\documentclass[10pt,a4paper]{amsart}
\usepackage[margin=19mm]{geometry}
\usepackage{amsmath,amssymb,mathtools}
\usepackage[scr=rsfs,cal=euler]{mathalfa}
\usepackage{xfrac}
\usepackage[unicode,hidelinks,bookmarks=false]{hyperref}
\newcommand{\ie}{i.e.\ }
\newcommand{\cf}{cf.\ }
\newcommand{\resp}{respectively\ }
\newcommand{\Subsection}{\subsection}
\providecommand{\nobreakdash}{\nobreak\hbox{-}}
\setlength{\emergencystretch}{2em}
\multlinegap0pt
\usepackage{bm}
\usepackage{bbm}
\usepackage{dsfont}
\usepackage{xspace}
\usepackage{cancel}
\usepackage{mathtools}
\usepackage{amsthm}
\makeatletter
\@ifundefined{equation}{\newtheorem{equation}{Equation}}{}
\@ifundefined{Proposition}{\newtheorem{Proposition}[equation]{Proposition}}{}
\@ifundefined{Theorem}{\newtheorem{Theorem}[equation]{Theorem}}{}
\makeatother
\renewcommand{\Re}[1]{\operatorname{Re} #1 }
\renewcommand{\geq}{\geqslant}
\renewcommand{\leq}{\leqslant}
\title[Crescents and the real variable Cesaro operator]{Crescents and the real variable Cesaro operator}
\author{Anil Belli, Ugur Gul, William T. Ross, Aristos Siskakis}
\date{Source: arXiv:2606.08600v1 (CC BY 4.0)}
\begin{document}
\maketitle

\noindent\emph{Source and licence.} Crescents and the real variable Cesaro operator. Version arXiv:2606.08600v1 (CC BY 4.0), distributed under the Creative Commons Attribution 4.0 International licence, as recorded in the arXiv licence metadata for this exact version and shown on the abstract page https://arxiv.org/abs/2606.08600v1. Full rights notice: licenses/arxiv-2606.08600-CC-BY-4.0.txt.\par\smallskip

\noindent\textit{Editorial scope.} Counted unit: the Theorem whose source label is \texttt{rrrrAAp} in the article (source0.tex lines 514--544, source label \texttt{rrrrAAp}), delivered together with the article's own complete original proof of it (source0.tex lines 546--622, one \texttt{proof} environment of 77 lines). No mathematics is written, repaired or re-derived: statement and proof are the article's own text line for line. Required context retained uncounted: 9oijrteglkfdsfawfgrdgfds1 (lines 286--294); www000 (lines 363--368); LPppp (lines 408--410). Every definition, display and label of those lines is retained unchanged. Omitted from this package and not redistributed: 1--285, 295--362, 369--407, 411--513, 545--545, 623--1002. Nothing inside a retained excerpt is omitted or altered: every retained body is delivered exactly as the article's lines are written, so no omission receipt of any kind accompanies this package. Citations: the retained excerpts carry no citation command at all, so no bibliography entry of the article is transcribed and no reference is redistributed. Reference adaptation: none was needed and none was made; every reference inside the delivered unit resolves inside it, so no unresolved reference or citation is produced. Printed slips kept literal: no printed word, symbol, hypothesis or number of the counted statement or of its proof was altered. Layout and class: the article's own class is \texttt{amsart} with packages including amsmath, amssymb, amsthm, amsfonts, enumerate, color, comment, mathrsfs, parskip, xcolor, mathtools, palatino, graphicx, enumitem; the wrapper is the lane's pinned \texttt{amsart} editorial wrapper at 10pt on A4, which restates the theorem-like environments the retained text needs and adds the article's own macro definitions that the retained text needs (\texttt{\textbackslash Re}, \texttt{\textbackslash geq}, \texttt{\textbackslash leq}), with extra wrapper packages (bm, bbm, dsfont, xspace, cancel, mathtools). The delivered numbering is the unit's own. Assets: no retained span contains a figure environment, a graphics-inclusion command, a PDF-page inclusion, a table or a TikZ picture; no image file is copied into this package, the package declares no asset dependency and no image file is redistributed with it. Licence: version arXiv:2606.08600v1 of this article is distributed under the Creative Commons Attribution 4.0 International licence, as recorded in the arXiv licence metadata for this exact version and shown on the abstract page https://arxiv.org/abs/2606.08600v1. A full rights notice is delivered at licenses/arxiv-2606.08600-CC-BY-4.0.txt. Characters: every retained excerpt is pure ASCII, so the delivered engine, XeLaTeX with its default Latin Modern text font, reports no missing character.
\begin{Proposition}\label{9oijrteglkfdsfawfgrdgfds1}
For the semigroup $\{S_t\}_{t \geq 0}$ we have the following. 
\begin{enumerate}
%\item If $p = 1$ then $e^{-t} \leq \|S_{t}\|_{L^1 \to L^1} \leq 1$ for all $t \geq 0$.
\item If $1 < p < 2$ then $e^{-\frac{t}{p}} \leq \|S_{t}\|_{L^p \to L^p} \leq e^{-\frac{t}{q}}$ for all $t \geq 0$. 
%\item If $p = 2$ then $\|S_{t}\|_{L^2 \to L^2} = e^{-\frac{t}{2}}$ for all $t \geq 0$.
\item If $2 \leq p < \infty$ then $e^{-\frac{t}{q}} \leq \|S_{t}\|_{L^p \to L^p} \leq e^{-\frac{t}{p}}$ for all $t \geq 0$. 
\end{enumerate}
\end{Proposition}

\begin{equation}\label{www000}
\omega_0 \leq \begin{cases}
-\frac{1}{q}& \mbox{if $1 < p < 2$,}\\[4pt]
-\frac{1}{p} &\mbox{if $p \geq 2$.}
\end{cases}
\end{equation}

\begin{equation}\label{LPppp}
R(\lambda,A_p)f =(\lambda I-A_p)^{-1}f =\int_{0}^{\infty}e^{-\lambda t}(S_{t}f)dt.
\end{equation}

\begin{Theorem}\label{rrrrAAp}
\hfill
\begin{enumerate}
\item For each $1 < p < 2$ we have the following.
\begin{enumerate}
\item Whenever $\Re \lambda> -\frac{1}{q}$,
\begin{equation}\label{aone}
 R(\lambda,A_p)f(x)=(\lambda I-A_p)^{-1}f(x)=\frac{(1-x)^{\lambda}}{x^{\lambda+1}}\int_{0}^{x}\frac{t^{\lambda}f(t)dt}{(1-t)^{\lambda+1}}.
 \end{equation}
\item Whenever $\Re \lambda<-\frac{1}{p}$,
\begin{equation}\label{atwo}
 R(\lambda,A_p)f(x)=(\lambda I-A_p)^{-1}f(x)=-\frac{(1-x)^{\lambda}}{x^{\lambda+1}}\int_{x}^{1}\frac{t^{\lambda}f(t)dt}{(1-t)^{\lambda+1}}.
 \end{equation}
  \end{enumerate}
 Moreover, $R(\lambda,A_p)$ is a bounded linear operator on $L^{p}(0,1)$ whenever 
$$ \Re \lambda > -\tfrac{1}{q} \; \mbox{or} \; \Re \lambda < -\tfrac{1}{p}.$$
\item For each  $p \geq 2$ we have the following.
\begin{enumerate}
\item Whenever $\Re \lambda > -\frac{1}{p}$, 
\begin{equation}\label{bone}
 R(\lambda,A_p)f(x)=(\lambda I-A_p)^{-1}f(x)=\frac{(1-x)^{\lambda}}{x^{\lambda+1}}\int_{0}^{x}\frac{t^{\lambda}f(t)dt}{(1-t)^{\lambda+1}}.
 \end{equation}
 \item Whenever $\Re \lambda < -\frac{1}{q}$, 
 \begin{equation}\label{btwo}
  R(\lambda,A_p)f(x)=(\lambda I-A_p)^{-1}f(x)=-\frac{(1-x)^{\lambda}}{x^{\lambda+1}}\int_{x}^{1}\frac{t^{\lambda}f(t)dt}{(1-t)^{\lambda+1}}.
  \end{equation}
\end{enumerate}
   Moreover, $R(\lambda,A_p)$ is a bounded linear operator on $L^{p}(0,1)$ whenever 
$$\Re \lambda > -\tfrac{1}{p} \; \mbox{or} \; \Re \lambda < -\tfrac{1}{q}.$$
\end{enumerate}
\end{Theorem}

\begin{proof}
From \eqref{www000}, along with the discussion preceding it, we know  that $R(\lambda, A_p)$ is bounded on $L^p(0, 1)$ whenever $\Re \lambda>-\frac{1}{q}$ (when $1 < p < 2$) or whenever $\Re \lambda > -\frac{1}{p}$ (when $p \geq 2)$.
Moreover, in both cases, we use the Laplace transform identity in  \eqref{LPppp} to see that 
\begin{align*}
R(\lambda,A_p)f(x) & =\int_{0}^{\infty}e^{-t\lambda}S_{t}f(x)dt\\
& =\int_{0}^{\infty}(e^{-t})^{\lambda}\frac{e^{-t}}{(e^{-t}-1)x+1}f\bigg(\frac{e^{-t}x}{(e^{-t}-1)x+1}\bigg)dt.
\end{align*}
With the substitution 
$$u=\frac{e^{-t}x}{(e^{-t}-1)x+1},$$ 
and checking
$$e^{-t} = \frac{u (1 - x)}{x (1 - u)} \;  \; \mbox{and} \; \; dt = - \frac{1}{u (1 - u)} du,$$
the above integral transforms into
$$
R(\lambda,A_p)f(x) = 
 \frac{(1-x)^{\lambda}}{x^{\lambda+1}}\int_{0}^{x}\frac{u^{\lambda}f(u)du}{(1-u)^{\lambda+1}},
$$
which proves \eqref{aone} and \eqref{bone}.



On the other hand, a formula for $R(\lambda,A_p)$ is found by solving
$$(\lambda I-A_p)f(x)=\lambda f(x)+(1-x)(xf(x))'=g(x)$$
for $f$ in terms of $g$. The integrating factor
$$\mu(x):=\frac{x^{\lambda +1}}{(1-x)^{\lambda}}$$
transforms this differential equation into 
$$(\mu(x)f(x))'=\Big(\frac{x^{\lambda +1}}{(1-x)^{\lambda}}f(x)\Big)'=\frac{x^{\lambda}}{(1-x)^{\lambda +1}}g(x).$$
Whenever $1 < p < 2$ and $\Re \lambda <-\frac{1}{p}$, observe that
$$\lim_{x \to 1^{-}} \frac{x^{\lambda +1}}{(1-x)^{\lambda}}f(x) = 0$$ and, for each $0 < x < 1$, the integral 
$$\int_{x}^{1}\frac{t^{\lambda}}{(1-t)^{\lambda +1}}g(t)dt$$
converges.
Thus,
$$(\lambda I-A_p)^{-1}f(x)=-\frac{(1-x)^{\lambda}}{x^{\lambda+1}}\int_{x}^{1}\frac{t^{\lambda}f(t)dt}{(1-t)^{\lambda+1}},
$$
which proves \eqref{atwo}. A similar argument verifies \eqref{btwo} when $p \geq 2$ and $\Re \lambda < -\frac{1}{q}$. 


It remains to prove $R(\lambda,A_p)$ is bounded on $L^{p}(0,1)$ whenever $\Re \lambda <-\frac{1}{p}$ (for $1 < p < 2$) or whenever $\Re \lambda < -\frac{1}{q}$ (for $p \geq 2$). To this end, let 
$$I_{\mu} f:=\int_{0}^{\infty}e^{\mu t}S_{t}fdt.$$
Formally substituting 
$$u=\frac{e^{-t}x}{(e^{-t}-1)x+1}$$
in the above integral and 
and checking that 
$$e^{t} = \frac{x(1 - u)}{u (1 - x)} \;  \; \mbox{and} \; \; dt = - \frac{1}{u (1 - u)} du,$$
 we   have 
\begin{align*}
I_{\mu}f(x)& =\int_{0}^{x}(e^{t})^{\mu}\frac{u}{x}\frac{f(u)}{u(1-u)}du\\ & =\int_{0}^{x}\bigg(\frac{x(1-u)}{u(1-x)}\bigg)^{\mu}\frac{1}{x}\frac{f(u)}{1-u}du\\
&=\frac{x^{\mu-1}}{(1-x)^{\mu}}\int_{0}^{x}\frac{(1-u)^{\mu-1}}{u^{\mu}}f(u)du.
\end{align*}
Now consider the formula 
$$R(\lambda,A_p)f(x)=(\lambda I-A_p)^{-1}f(x)=-\frac{(1-x)^{\lambda}}{x^{\lambda+1}}\int_{x}^{1}\frac{t^{\lambda}f(t)dt}{(1-t)^{\lambda+1}},$$
from (b). 
Make the substitution $u=1-t$ to get 
$$R(\lambda,A_p)f(x)=-\frac{(1-x)^{\lambda}}{x^{\lambda+1}}\int_{0}^{1-x}\frac{(1-u)^{\lambda}}{u^{\lambda+1}}f(1-u)du.$$
The  ``flip operator''
$(\mathcal{F}f)(x):=f(1-x)$
is isometric on $L^p(0, 1)$ and 
$$R(\lambda,A_p)=-\mathcal{F}\circ I_{\lambda+1}\circ\mathcal{F}.$$
Thus  on $L^p(0, 1)$, $R(\lambda,A_p)$ is bounded if and only if  $I_{\lambda+1}$ is bounded. We also have 
\begin{align*}
\| I_{\lambda+1}f\|_{L^p \to L^p}  & \leq\bigg(\int_{0}^{\infty}e^{(\Re (\lambda)+1)t}\| S_{t}f\|_{p}dt\bigg)\\
 & \leq\bigg(\int_{0}^{\infty}e^{(\Re(\lambda)+1-\frac{1}{q})t}dt\bigg)  \|f\|_p
 \end{align*}
(note the use of Proposition \ref{9oijrteglkfdsfawfgrdgfds1}(a)). Hence, since  
$$\int_{0}^{\infty}e^{(\Re(\lambda)+1-\frac{1}{q})t}dt<\infty$$ whenever 
$$\Re \lambda <-1+\tfrac{1}{q}=-(1-\tfrac{1}{q})=-\tfrac{1}{p},$$ $I_{\lambda+1}$ is bounded on $L^{p}(0, 1)$ whenever $1 < p < 2$ and $\Re \lambda<-\frac{1}{p}$. Conclusion: $R(\lambda,A_p)$ is bounded on $L^{p}(0, 1)$. 


In a similar way, 
\begin{align*}
\| I_{\lambda+1}f\|_{L^p \to L^p}  & \leq\bigg(\int_{0}^{\infty}e^{(\Re (\lambda)+1)t}\| S_{t}f\|_{p}dt\bigg)\\
 & \leq\bigg(\int_{0}^{\infty}e^{(\Re(\lambda)+1-\frac{1}{p})t}dt\bigg)  \|f\|_p,
 \end{align*}
 (note the use of Proposition \ref{9oijrteglkfdsfawfgrdgfds1}(b)) 
 and since 
$$\int_{0}^{\infty}e^{(\Re(\lambda)+1-\frac{1}{p})t}dt<\infty$$ whenever 
$$\Re \lambda <-1+\tfrac{1}{p}=-(1-\tfrac{1}{p})=-\tfrac{1}{q},$$  $I_{\lambda+1}$ is bounded on $L^{p}(0, 1)$ whenever $p \geq 2$ and $\Re \lambda<-\frac{1}{q}$. As before, this implies $R(\lambda,A_p)$ is bounded on $L^{p}(0, 1)$. 
\end{proof}

\end{document}
